\documentclass[10pt,a4paper]{amsart}
\usepackage[margin=19mm]{geometry}
\usepackage{amsmath,amssymb,mathtools}
\usepackage[scr=rsfs,cal=euler]{mathalfa}
\usepackage{xfrac}
\usepackage[unicode,hidelinks,bookmarks=false]{hyperref}
\newcommand{\ie}{i.e.\ }
\newcommand{\cf}{cf.\ }
\newcommand{\resp}{respectively\ }
\newcommand{\Subsection}{\subsection}
\providecommand{\nobreakdash}{\nobreak\hbox{-}}
\setlength{\emergencystretch}{2em}
\multlinegap0pt
\usepackage{amssymb}
\usepackage{tikz}
\usepackage{float}
\usepackage{xcolor}
\newtheorem{lem}{Lemma}
\title{Fluid limit of a Distributed Ledger Model with Random Delay}
\author{Jiewei Feng, Christopher King}
\date{Source: arXiv:2311.03475v2 (CC BY 4.0)}
\begin{document}
\maketitle

\noindent\emph{Source and licence.} Fluid limit of a Distributed Ledger Model with Random Delay. Version arXiv:2311.03475v2 (CC BY 4.0), distributed by arXiv under the Creative Commons Attribution 4.0 International licence, http://creativecommons.org/licenses/by/4.0/, as recorded in the arXiv licence metadata for this exact version and shown on the abstract page https://arxiv.org/abs/2311.03475v2. Full licence notice: \texttt{licenses/arxiv-2311.03475-CC-BY-4.0.txt}.\par\smallskip

\noindent\textit{Editorial scope.} Counted unit: the article's lem k\_extendw (source0.tex lines 841--848). The article's complete original proof of it is delivered with it (source0.tex lines 1437--1599, one \texttt{proof} environment of 163 lines, which the article places away from the statement). No mathematics is written, repaired or re-derived: the statement and the proof are the article's own lines, in the article's own notation, and no retained line is substituted or re-flowed.  No further source-local context is needed: every cross-reference of every retained line resolves inside the two delivered spans.  Nothing was omitted from any retained span, so the package carries no omission record.  No retained line contains a citation, so the wrapper carries no bibliography and no bibliography entry.  No retained line contains a figure environment, an image-inclusion command or drawing code, so no image is delivered and the package declares no asset dependency.  Editorial formatting only, in addition: an AMS \texttt{amsart} preamble that restates the article's own declarations and macros used by the retained lines, namely the environments \texttt{lem} on separate counters, together with the standard packages those lines need.  The retained environments print in this document as Lemma 1 in order of appearance, whereas the article prints the same items with the numbering of its own full document.  The complete native source of the article as distributed in the arXiv e-print for this exact version is bundled unchanged as \texttt{sources/149623/source0.tex}.
\begin{lem}\label{k_extendw}%{restatable}{lemma}{temptwelve}\label{extendw}
Consider $s_1,s_2\in[k\epsilon,(k+1)\epsilon)$ and $s_2>s_1$. Then for $i=1,2,...,M$, $\left|w_i(t,s_2)-w_i(t,s_1)\right|\leq \alpha_{0,1}M(h_M)^M\exp(2Mh_M/m)\epsilon$, where 
\begin{align*} 
\alpha_{0,1}=(2+2Mh_M)\xi+6/m+
+\frac{4(3+4M h_M \xi)}{m}
+2h_M\xi \frac{2(3+4 M h_M \xi)}{m}.
\end{align*} 
\end{lem}

\begin{proof}[Proof of Lemma \ref{k_extendw}]
It is easier for proof to replace $s_j$ by $h_i-s_j$, that is, we focus on
\begin{align*} 
w_i(t,h_i-s_1)&= w_i\big(t-s_1,h_i\big)+\int_{0}^{s_1} \sum_{h_j\leq h_i-u} -2p_j\frac{w_i\big(t-s_1+u,h_i-u\big)}{l(t-s_1+u)}du,
\end{align*} 
and
{
\begin{align*} 
&w_i(t,h_i-s_2)= w_i\big(t-s_2,h_i\big)+\int_{0}^{s_2} \sum_{h_j\leq h_i-u} -2p_j\frac{w_i\big(t-s_2+u,h_i-u\big)}{l\big(t-s_2+u\big)}du\\
&=w_i(t-s_2,h_i)+\int_{s_1-s_2}^{s_1} \sum_{h_j\leq h_i-(u-s_1+s_2)} -2p_j\frac{w_i\big(t-s_1+u,h_i-\big(u-s_1+s_2\big)\big)}{l\big(t-s_1+u\big)} du.
\end{align*} }

Hence we have the following
{
\begin{align} 
&\left|w_i(t,h_i-s_2)-w_i(t,h_i-s_1)\right|\nonumber\\
\leq& \left|w_i\big(t-s_2,h_i\big)- w_i\big(t-s_1,h_i\big)\right|
+\left|\int_{s_1-s_2}^0-2\sum_{h_j} p_j\frac{w_i\big(t-s_1+u,h_i-\big(u-s_1+s_2\big)\big)}{l\big(t-s_1+u\big)}du\right|\nonumber\\
&+\int_{0}^{s_1}2\sum_{h_j\leq h_i-u}p_j \left|\frac{w_i\big(t-s_1+u,h_i-u\big)-w_i\big(t-s_1+u,h_i-\big(u-s_1+s_2\big)\big)}{l\big(t-s_1+u\big)}\right|du\nonumber\\
&+2\int_{0}^{s_1} 
\left|
\sum_{h_j\leq h_i-(u-s_1+s_2)} \frac{p_j w_i\big(t-s_1+u,h_i-\big(u-s_1+s_2\big)\big)}{l\big(t-s_1+u\big)} \right.\nonumber\\
&\quad\quad\quad\quad\quad\quad\quad\left.-\sum_{h_j\leq h_i-u} \frac{p_j w_i\big(t-s_1+u,h_i-\big(u-s_1+s_2\big)\big)}{l\big(t-s_1+u\big)} 
\right|du .
\label{kpw_12_0}
\end{align} }


For the first term on the RHS of (\ref{kpw_12_0}),
\begin{align}
&\left|w_i\big(t-s_2,h_i\big)- w_i\big(t-s_1,h_i\big)\right|\nonumber\\
&\leq\left|2p_i\frac{f\big(t-s_2\big)}{l\big(t-s_2\big)}-2p_i\frac{f\big(t-s_1\big)}{l\big(t-s_1\big)}\right|
+\sum_{j>i}\int_{ h_i}^{h_j}2p_i
\left|
\frac{w_j\big(t-s_1,u\big)}{l\big(t-s_1\big)}-\frac{w_j\big(t-s_2,u\big)}{l\big(t-s_2\big)}
\right| du
\nonumber\\
&\leq2\left[\frac{f\big(t-s_2\big)-f\big(t-s_1\big)}{l\big(t-s_2\big)}+\frac{f\big(t-s_1\big)}{l\big(t-s_1\big)}\frac{l\big(t-s_1\big)-l\big(t-s_2\big)}{l\big(t-s_2\big)}\right]+..\nonumber.\\
&\leq \frac{6\epsilon}{m}+\frac{2(3+4M h_M \xi)\epsilon}{m}
+\sum_{j>i}\int_{ h_i}^{h_j}2p_i
\left|
\frac{w_j\big(t-s_1,u\big)}{l\big(t-s_1\big)}-\frac{w_j\big(t-s_2,u\big)}{l\big(t-s_2\big)}
\right| du\nonumber\\
&\leq\frac{6\epsilon}{m}+\frac{2\big(3+4M h_M \xi\big)\epsilon}{m}+2h_M\xi \frac{2\big(3+4M h_M\xi\big)\epsilon}{m}\nonumber\\
&\quad+\sum_{j>i}\int_{ h_i}^{h_j}2p_i
\left|
w_j\big(t-s_1,u\big)-w_j\big(t-s_2,u\big)
\right|
du\label{kpw_12_1},
\end{align}
where we use the facts $f(t)<l(t)$, $s_2-s_1<\epsilon$, $l(t)>m$, and $\xi>w_j$. For the second term on the RHS of (\ref{kpw_12_0}),
\begin{align}
\left|\int_{s_1-s_2}^0-2 \sum_{h_j}p_j \frac{w_i\big(t-s_1+u,h_i-\big(u-s_1+s_2\big)\big)}{l\big(t-s_1+u\big)}du\right|\leq\int_{0}^{\epsilon}2\xi
\leq 2\xi\epsilon\label{kpw_12_2},
\end{align}
where $\xi$ is an upper bound for $\frac{w_2(s+u,u)}{l(s+u)}$ for all $s,u$. For the last term on the RHS of (\ref{kpw_12_0}), because $|s_1-s_2|<\epsilon$, we have
{
\begin{align} 
&2\int_{0}^{s_1} 
\left|
\sum_{h_j\leq h_i-(u-s_1+s_2)} p_j\frac{w_2\big(t-s_1+u,h_i-\big(u-s_1+s_2\big)\big)}{l\big(t-s_1+u\big)}\right.\nonumber\\
&\left.\quad\quad\quad\quad\quad
-\sum_{h_j\leq h_i-u} p_j\frac{w_2\big(t-s_1+u,h_i-\big(u-s_1+s_2\big)\big)}{l\big(t-s_1+u\big)} 
\right|
du\nonumber\\
&\leq 2\int_{0}^{s_1} 
\left|
\sum_{h_i-u <h_j\leq h_i-\big(u-s_1+s_2\big)} p_j\frac{w_2\big(t-s_1+u,h_i-\big(u-s_1+s_2\big)\big)}{l\big(t-s_1+u\big)} 
\right| 
du\nonumber\\
&\leq 2\sum_{j<i}\int_{h_i-h_j+s_2-s_i}^{h_i-h_j}
\left|
p_j\frac{w_2\big(t-s_1+u,h_i-\big(u-s_1+s_2\big)\big)}{l\big(t-s_1+u\big)} 
\right| 
du\leq 2 \xi \epsilon \label{kpw_12_4}.
\end{align}}

All that left is the third term in (\ref{kpw_12_0}). Define
\begin{align}
y_i(u)=\left|w_i\big(t-s_1+u,h_i-u\big)-w_i\big(t-s_1+u,h_i-\big(u-s_1+s_2\big)\big)\right|\label{kpw_12_5}.
\end{align}

Then we have $\left|w_i(t,h_i-s_2)-w_i(t,h_i-s_1)\right|=y_i(s_1)$ and substituting (\ref{kpw_12_1}), (\ref{kpw_12_2}), (\ref{kpw_12_4}), and (\ref{kpw_12_5}) into (\ref{kpw_12_0}) we get
\begin{align*} 
y_i(s_1)&\leq 4\xi\epsilon+\frac{6\epsilon}{m}+\frac{2(3+4M h_M\xi)\epsilon}{m}
+2h_M\xi \frac{2(3+4M h_M \xi)\epsilon}{m}\\
&+\sum_{j>i}\int_{ h_i}^{h_j}2p_i
\left|
w_j\big(t-s_1,u\big)-w_j\big(t-s_2,u\big)
\right|
du+\frac{2}{m}\int_{0}^{s_1}y_i(u)du.
\end{align*} 
And since
{
\begin{align*} 
&\left|
w_j\big(t-s_1,u\big)-w_j\big(t-s_2,u\big)
\right|\\
\leq& 
\left|
w_j\big(t-s_1,u\big)-w_j\big(t-s_1,u-s_2+s_1\big)
\right|\nonumber\\
&+2\left|
\int_{0}^{s_1-s_2}\sum_{h_r\leq u-s_2+s_1+s}p_j\frac{w_j\big(t-s_1+s,u-s_2+s_1+s\big)}{l\big(t-s_2+s\big)}du
\right|\nonumber\\
\leq&\left|
w_j\big(t-s_1,u\big)-w_j\big(t-s_1,u-s_2+s_1\big)
\right|+2\xi\epsilon,
\end{align*}}
we get
\begin{align} 
&y_i(s_1)\leq 4\xi\epsilon+\frac{6\epsilon}{m}+\frac{2(3+4M h_M \xi)\epsilon}{m}
+2h_M\xi \frac{2(3+4M h_M \xi)\epsilon}{m}\nonumber\\
&+2Mh_M\xi\epsilon+\sum_{j>i}\int_{ h_i}^{h_j}2p_i
\left|
w_j\big(t-s_1,u\big)-w_j\big(t-s_1,u-s_2+s_1\big)
\right|
du+\frac{2}{m}\int_{0}^{s_1}y_i(u)du\label{kpw_12_6}.
\end{align} 
We now use induction on $i$ to (\ref{kpw_12_6}).

\textbf{Case 1} : $i=M$. There is no term in the summation $\sum_{j>i}$ included in (\ref{kpw_12_6}). So 
\begin{align*} 
y_i\big(s_1\big)\leq& (2+2Mh_M)\xi\epsilon+\frac{6\epsilon}{m}+\frac{2(3+4 M h_M \xi)\epsilon}{m}
+\frac{2(3+4M h_M \xi)\epsilon}{m}\\
&+2h_M\xi \frac{2(3+4 M h_M \xi)\epsilon}{m}
+\frac{2}{m}\int_{0}^{s_1}y(u)du\\
=:&\alpha_{1,0}\epsilon+\frac{2}{m}\int_{0}^{s_1}y_i(u)du
\end{align*} 
Applying Gronwall's inequality we get
\begin{align*} 
y_M(s_i)=\left|w_M(t,h_M-s_2)-w_M(t,h_M-s_1)\right|\leq\alpha_{1,0}\epsilon \exp\left\{\int_{0}^{s_1}\frac{2}{m}du\right\}\leq\alpha_{1,0}\epsilon \exp\left\{\frac{2h_M}{m}\right\}.
\end{align*} 

\textbf{Case 2}: $i=M-1$.
Using the result for case 1, we get
\begin{align*} 
y_i(s_1)\leq& \alpha_{1,0}\epsilon
+\sum_{j>i}\int_{ h_i}^{h_j}2p_i
\left|
w_j\big(t-s_1,u\big)-w_j\big(t-s_2,u\big)
\right|
du+\frac{2}{m}\int_{0}^{s_1}y_i(u)du\\
\leq& \alpha_{1,0}\epsilon
+\int_{ h_i}^{h_M}2
\left|
w_M\big(t-s_1,u\big)-w_M\big(t-s_2,u\big)
\right|
du +\frac{2}{m}\int_{0}^{s_1}y_i(u)du\\
\leq& \alpha_{1,0}\epsilon+2h_M\alpha_{1,0} \exp\left\{\frac{2h_M}{m}\right\}\epsilon+\frac{2}{m}\int_{0}^{s_1}y_i(u)du\\
\leq&(3h_M)\alpha_{1,0} \exp\left\{\frac{2h_M}{m}\right\}\epsilon+\frac{2}{m}\int_{0}^{s_1}y_i(u)du.
\end{align*} 

Again applying Gronwall's inequality we get
\begin{align*} 
\left|w_{M-1}(t,s_2)-w_{M-1}(t,s_1)\right|\leq(2h_M+1)\alpha_{1,0} \exp\left\{\frac{4h_M}{m}\right\}\epsilon. 
\end{align*} 
Repeat and iterate on $i$ we get for all $i\leq M$
we get 
\begin{align*} 
\left|w_i(t,s_2)-w_i(t,s_1)\right|\leq \alpha_{0,1}M\big(h_M\big)^M\exp\left\{\frac{2Mh_M}{m}\right\}\epsilon. 
\end{align*} 
\end{proof}

\end{document}
